\section{Guarantees and Correctness Properties}
\label{sec:correctness-properties}

The preceding sections establish the Bailout Protocol as a formal mechanism: a deterministic state machine, a trigger condition, an escalation algorithm, and a bypass mechanism enforcing the no-machine-actor property. This section formally states the three correctness properties the protocol is designed to guarantee, identifies the structural assumptions on which each guarantee rests, and provides a proof sketch for each. All three properties are stated as invariants over the full state space of $\mathcal{B}$, not merely for a specific execution path.

% ---------------------------------------------------------------
\subsection{Property 1: Safety — No Autonomous Collapse}
% ---------------------------------------------------------------

\theorem[Safety (No Autonomous Collapse)]{During any execution of the Bailout Protocol, the system shall not transition to a terminal resolving state without an explicit, recorded human authorization.}

\subparagraph{Formal statement.} Let $\mathcal{B}$ be the Bailout Protocol state machine for node $n$. Let $P(t)$ denote the set of functional communication pathways to $h(n)$ available at time $t$. Safety requires:
\begin{equation}
\forall t : |P(t)| = 0 \;\land\; \text{state}(n) = \text{RESOLVED}
\;\implies\;
\text{tag}(n) = \texttt{unresolved}
\end{equation}
Equivalently: if no pathway to the human anchor is available, the system may enter RESOLVED only with an \texttt{unresolved} tag — a documented stall, not an autonomous resolution.

\subparagraph{Structural assumptions.}
\begin{enumerate}
\item[\textbf{A1.1}] The Fact Layer pre-commit hook is operative: any \texttt{UpdateLedger} call is processed atomically and rejects any write that would record a directive issued by a machine actor.
\item[\textbf{A1.2}] The immutable parent pointer $\alpha(n)$ and human-root identifier $h(n)$ are established at node provisioning and are not writable by any machine actor at runtime.
\item[\textbf{A1.3}] At least one functional communication pathway to $h(n)$ is provisioned at node startup. This is a deployment requirement, not a protocol invariant.
\end{enumerate}

\subparagraph{Proof sketch.} The proof proceeds by enumerating all transitions into the terminal state.

\begin{itemize}
\item[Step 1.] The only transitions entering RESOLVED are T7 (ESCALATING $\rightarrow$ RESOLVED via human RESOLVE or REJECT) and T9 (HUMAN\_ACKNOWLEDGED $\rightarrow$ RESOLVED via human RESOLVE or REJECT). Both require a directive $d \in \{\texttt{RESOLVE}, \texttt{REJECT}\}$ issued by $h(n)$. By INV\_2, the Fact Layer enforces $\text{directive\_issuer}(n) = h(n)$ at the authorization ledger level. Any attempted entry to RESOLVED with a machine-issued directive is rejected by the pre-commit hook; the machine stalls in its current non-terminal state.

\item[Step 2.] Consider the two source states for RESOLVED:
\begin{enumerate}
\item[\textbf{ESCALATING}:] The \texttt{Escalate} algorithm propagates the exception along the parent chain using the Pathway Health Registry. If all provisioned pathways fail $N_{\max}$ consecutive health checks, the algorithm fires a \texttt{parent\_unreachable} timeout event, propagates it to $h(n)$ via the next available functional pathway, and records the failure in the Forensic Ledger. If $|P(t)| = 0$, no propagation can succeed; the originating node stalls in ESCALATING. No machine actor can issue a directive from this state by BP (bypass property). Therefore RESOLVED cannot be entered without human authorization when all pathways are unavailable.
\item[\textbf{HUMAN\_ACKNOWLEDGED}:] Entry occurs only via T2 or T6, both requiring a Fact-Layer-recorded human acknowledgment. Only $h(n)$ may issue a directive from this state. If the human anchor does not respond before $T_{\text{human}}$ expires, the protocol transitions to RESOLVED with an \texttt{unresolved} tag — a terminal state that documents a stall condition, not an autonomous resolution.
\end{enumerate}
\item[Step 3.] The trigger condition TC and the bypass mechanism BP jointly prevent any state in $Q \setminus \{\text{RESOLVED}\}$ from being reached via a machine-actor-issued directive. Reaching RESOLVED without authorization would require the Fact Layer to approve a non-human directive, violating A1.1. Safety thus holds conditionally on A1.1.
\end{itemize}

\subparagraph{Discussion.} Safety does not require a pathway to $h(n)$ to always exist — it requires that the system cannot autonomously collapse regardless of pathway availability. When all pathways fail, the system stalls in ESCALATING or enters RESOLVED with \texttt{unresolved}. Both outcomes are documented stalls, not unauthorized resolutions. The distinction is operationally meaningful: in safety-critical domains, a system that fails by stalling and logging the failure is preferable to one that proceeds on an unauthorized resolution.

% ---------------------------------------------------------------
\subsection{Property 2: Liveness — Eventual Human Contact}
% ---------------------------------------------------------------

\theorem[Liveness (Eventual Human Contact)]{During any execution of the Bailout Protocol, if the system enters BOUNDED and the Bounds Engine does not produce a Fact-Layer-approved resolution within $T_{\text{bounds}}$, then the protocol must eventually enter either HUMAN\_ACKNOWLEDGED or RESOLVED (with an \texttt{unresolved} tag), and in either case $h(n)$ is contacted within a bounded amount of wall-clock time.}

\subparagraph{Formal statement.} Let $t_0$ be the time at which T1 fires (entry to BOUNDED). Liveness requires:
\begin{equation}
\exists t_1 \geq t_0 : t_1 - t_0 \leq T_{\max}
\;\implies\;
\text{state}(n) \in \{\text{HUMAN\_ACKNOWLEDGED}, \text{RESOLVED}\}
\end{equation}
where $T_{\max}$ is a deployment-specific bound. If $h(n)$ does not respond within $T_{\text{human}}$, the protocol enters RESOLVED with an \texttt{unresolved} tag; this constitutes human contact because the notification was sent and received, as evidenced by the Forensic Ledger entry.

\subparagraph{Structural assumptions.}
\begin{enumerate}
\item[\textbf{A2.1}] The Pathway Health Registry is operative: at least one functional pathway to $h(n)$ is available at trigger time, or becomes available before $N_{\max}$ retries are exhausted.
\item[\textbf{A2.2}] The human anchor $h(n)$ is responsive: absent a directive within $T_{\text{human}}$, the timeout issues a reminder and transitions to RESOLVED with \texttt{unresolved}, recorded in the Forensic Ledger.
\item[\textbf{A2.3}] Timeout parameters $T_{\text{bounds}}$, $T_{\text{prop}}$, $T_{\text{cap}}$, and $T_{\text{human}}$ are provisioned with finite positive values and are not adjustable by any machine actor at runtime.
\end{enumerate}

\subparagraph{Proof sketch.} The proof establishes that the protocol always advances from BOUNDED to a terminal state within a bounded time.

\begin{itemize}
\item[Step 1.] From BOUNDED, exactly two exits exist: T2 (HUMAN\_ACKNOWLEDGED via intra-bounds resolution) or T3 (BAIL\_PENDING $\rightarrow$ ESCALATING upon $T_{\text{bounds}}$ expiration or explicit \texttt{bailout\_signal}). T2 is the fast path; T3 is the guaranteed path. Since $T_{\text{bounds}}$ is a finite provisioned constant (A2.3), T3 fires within $T_{\text{bounds}}$ if T2 has not fired. The Fact Layer enforces $T_{\text{bounds}}$ at the pre-commit hook; no machine actor can prevent its expiration.

\item[Step 2.] T4 (BAIL\_PENDING $\rightarrow$ ESCALATING) fires immediately upon entry to BAIL\_PENDING with no dwell time — an architectural compulsion, not a configurable delay. From the moment T3 fires, the protocol enters ESCALATING within a negligible wall-clock delta bounded by atomic state transition latency.

\item[Step 3.] In ESCALATING, the \texttt{Escalate} algorithm iterates: select the lowest-latency functional pathway from the PHR, send the exception to the parent node, retry with exponential backoff capped at $T_{\text{cap}}$ if no acknowledgment arrives within $T_{\text{prop}}$. The maximum retry count $N_{\max}$ is a deployment constant. Total propagation time from ESCALATING entry to parent acknowledgment or $N_{\max}$-retry exhaustion is bounded by $N_{\max} \cdot T_{\text{cap}}$ plus accumulated backoff.

\item[Step 4.] If the parent is unreachable after $N_{\max}$ retries, the algorithm fires \texttt{parent\_unreachable} (T3 from ESCALATING), propagating to $h(n)$ via the next available functional pathway in the PHR. The fallback pathway selection prevents deadlock on a single failed parent; propagation advances to $h(n)$ via an alternate route.

\item[Step 5.] Upon reaching $h(n)$, the protocol enters HUMAN\_ACKNOWLEDGED (T6). If $h(n)$ does not issue a directive before $T_{\text{human}}$ expires (A2.2), the protocol transitions to RESOLVED with \texttt{unresolved} tag (T10). Both outcomes constitute human contact as defined by the liveness property.
\end{itemize}

\subparagraph{Combined bound.} The total wall-clock time from $t_0$ to human contact is bounded by:
\begin{equation}
T_{\max} \leq T_{\text{bounds}} + N_{\max} \cdot T_{\text{cap}} + T_{\text{human}}
\end{equation}
This bound is a deployment parameter. A deployment that cannot satisfy its operational $T_{\max}$ requirement must provision additional pathways, increase $N_{\max}$, or reduce timeout values — but the protocol guarantees human contact within whatever bound the deployment provisions.

\subparagraph{Discussion.} Liveness is conditional: it holds if A2.1, A2.2, and A2.3 are satisfied. If the human anchor is permanently unreachable, the protocol stalls in ESCALATING and records the failure. Safety (Property 1) ensures this stall does not become an autonomous resolution. Together, Safety and Liveness define the protocol's failure semantics: when human contact cannot be achieved because all pathways have failed, Safety guarantees the system fails by stalling, never by proceeding without authorization.

% ---------------------------------------------------------------
\subsection{Property 3: Accountability — Human Always Identifiable}
% ---------------------------------------------------------------

\theorem[Accountability (Human Always Identifiable)]{For every protocol run that produces a directive $d \in \{\texttt{ACK}, \texttt{RESOLVE}, \texttt{REJECT}\}$, the Forensic Ledger entry recording $d$ identifies the human principal who issued it, and no ledger entry records a directive issued by a machine actor.}

\subparagraph{Formal statement.} Let $\mathcal{L}$ be the append-only Forensic Ledger. Accountability requires:
\begin{equation}
\forall \ell \in \mathcal{L} : \ell.\text{event\_type} = \text{human\_directive}
\;\implies\;
\ell.\text{principal\_id} = h(n) \;\land\; \text{is\_human}(\ell.\text{principal\_id})
\end{equation}
where $h(n)$ is the human accountability anchor of the originating node $n$, and $\text{is\_human}$ returns true only for identifiers classified as human principal identifiers by the provisioning authority. Equivalently: every directive in the ledger is attributable to a human principal, and no machine-issued directive is recorded as valid.

\subparagraph{Structural assumptions.}
\begin{enumerate}
\item[\textbf{A3.1}] The provisioning authority issues human principal identifiers only to verified natural persons or institutional actors acting with unambiguous human authority. The $\text{is\_human}$ predicate is a deployment-side classification, not a protocol mechanism.
\item[\textbf{A3.2}] The Fact Layer pre-commit hook is operative and rejects any \texttt{UpdateLedger} call recording a directive whose $\text{principal\_id}$ field does not satisfy $\text{is\_human}$.
\item[\textbf{A3.3}] The parent pointer $\alpha(n)$ and human-root identifier $h(n)$ are established at provisioning and are immutable for the node's lifetime. Any runtime modification attempt is rejected by the Fact Layer pre-commit hook and logged as a protocol violation.
\end{enumerate}

\subparagraph{Proof sketch.} The proof proceeds by structural induction over protocol transitions that generate ledger entries.

\begin{itemize}
\item[Step 1.] The only transitions generating a ledger entry with $\text{event\_type} = \text{human\_directive}$ are T6 (ESCALATING $\rightarrow$ HUMAN\_ACKNOWLEDGED upon human ACK), T7 (ESCALATING $\rightarrow$ RESOLVED upon human RESOLVE or REJECT), T9 (HUMAN\_ACKNOWLEDGED $\rightarrow$ RESOLVED upon human RESOLVE or REJECT), and T10 (HUMAN\_ACKNOWLEDGED $\rightarrow$ RESOLVED upon human timeout with \texttt{unresolved} tag). Each fires only upon receipt of a directive from $h(n)$, as enforced by the Fact Layer pre-commit hook.

\item[Step 2.] The Fact Layer pre-commit hook enforces INV\_2: $\text{state}(n) = \text{HUMAN\_ACKNOWLEDGED} \implies \text{directive\_issuer}(n) = h(n)$. When a directive is received, the Fact Layer verifies the principal identifier against the authorization ledger. If the identifier does not satisfy $\text{is\_human}$ or does not match $h(n)$ for the originating node, the pre-commit hook rejects the directive and logs a protocol violation. This check cannot be bypassed by any machine actor, including the Purpose Layer: the Purpose Layer's role is exclusively receiving exceptions and issuing directives; it cannot authorize ledger entries attributing a directive to a non-human actor.

\item[Step 3.] The append-only property of the Forensic Ledger ensures that once written, no ledger entry can be modified or deleted. Even if a machine actor produced a directive with a falsified principal identifier, the Fact Layer's pre-commit check prevents the entry from being written. Falsification is blocked at authorization, not detected post-hoc.

\item[Step 4.] The immutability of $h(n)$ (A3.3) prevents a compromised node from substituting its own human anchor for the correct one. The $h(n)$ field is stored in Fact Layer-managed memory, read-only at the machine-actor level; any modification attempt is rejected and logged as an unauthorized modification.
\end{itemize}

\subparagraph{Discussion.} The accountability property guarantees attribution, not correctness: it guarantees that every directive in the ledger is attributable to a human principal, but not that the human principal made the correct decision. The latter is a separate concern — of human factors, organizational process, and training — that the protocol cannot address architecturally. What the protocol guarantees architecturally is that no machine actor can issue, simulate, or substitute a directive without this being detectable as a protocol violation. The ledger provides the documentary evidence; the pre-commit hook provides the enforcement. Together, they make accountability non-repudiable: a human principal cannot later deny issuing a directive that appears in the ledger under their identifier, and a machine actor cannot cause a directive to appear under a falsified human identifier.

The three properties are mutually reinforcing. Safety prevents the system from entering a terminal state without human authorization. Liveness ensures the system reaches human contact within a bounded time if a pathway exists. Accountability ensures that when human contact occurs, the resulting directive is attributable to a human and not to any machine actor. Safety and Liveness together define the protocol's failure semantics: when human contact cannot be achieved because all pathways have failed, Safety guarantees the system stalls rather than collapses autonomously, and the Forensic Ledger records the stall as a documented protocol event. Accountability ensures that even this failure mode is auditable — the system failure is recorded with full attribution to the human anchor who would have been the intended recipient.
